\chapter*{Glossary}
\addcontentsline{toc}{chapter}{Glossary}
Terms introduced in the first two chapters, collected here for reference.

\section*{Information and coding}
\begin{multicols}{2}\small\raggedright
\begin{rs}
\gloss{Pixel}: a location in an image grid, or its stored value.
\gloss{Symbol}: an element of the finite alphabet to be encoded.
\gloss{Codeword}: the bit string assigned to one symbol.
\gloss{Code}: the assignment of codewords to symbols.
\gloss{Empirical distribution}: the observed relative frequencies of symbols.
\gloss{Histogram}: a graphical representation of these frequencies.
\gloss{Source model}: a probability law for the symbols or their sequences.
\gloss{Entropy}: the nonnegative average information $H(V)$, in bits.
\gloss{Average code length}: $L(c)=\sum_vp_v\ell_v$, in bits per symbol.
\end{rs}
\end{multicols}

\section*{Image processing}
\begin{multicols}{2}\small\raggedright
\begin{rs}
\gloss{Bit and byte}: a binary digit; a group of eight bits.
\gloss{Channel}: one component array of a multichannel image.
\gloss{Quantization}: replacing values by indices from a finite set.
\gloss{Noise}: unwanted fluctuations introduced by acquisition or processing.
\gloss{Mean}: the sum of values divided by their number.
\gloss{Median}: the central sorted value for an odd-sized set.
\gloss{Edge}: a rapid spatial change in image values.
\gloss{Blur}: a loss of sharp transitions and fine detail.
\gloss{Contrast}: differences between light and dark regions.
\gloss{Luminance}: a photometric quantity defined using visual sensitivity.
\gloss{Matrix}: a rectangular array $(a_{i,j})$.
\gloss{Resolution}: a term that may refer to pixel dimensions or to spatial resolving power; these are distinct notions.
\gloss{Crossfade}: a weighted average of two images at fixed pixel locations.
\end{rs}
\end{multicols}
